\sscp{Differing terminological systems in sibling terms}{07-02-03}
We not only gain insights by looking at forms,
functions, and distributions of individual words or glosses,
but we also gain insights by looking at sets of related terms.
\citet{Fox2015} discusses the significance of different terminological
sets of kin terms for Austronesian languages.
Throughout our target region,
cognate kin terms may be used with a variety of different social structures.
The terminology does not determine the social structure.

In geographical central Maluku,
as throughout Wallacea, insights are to be gained by focusing
particularly on sets of sibling terms.
We see the following common patterns.

\begin{itemize}
	\item 2 terms (relative age): elder sibling :: younger sibling
	\item 4 terms: elder sibling (same sex) :: younger sibling (same sex) :: sister (male speaking) :: brother (female speaking)
	\item 3 terms: elder sibling (same sex) :: younger sibling (same sex) :: opposite sex sibling
	\item 2 terms (relative gender): opposite sex sibling :: same sex sibling
\end{itemize}

The two relative age sibling terms found widely in Sumatran,
\tsc{Celebic}, and Philippine languages are found with narrowed
meaning \ili{east} of the Blust line to add ‘same sex sibling’.
Consequently one or two ‘opposite sex sibling’ terms appear sporadically
in several subgroups \ili{east} of the Blust line.
While four reconstructed sibling terms have been assigned to
PMP and are referred to as retentions,
the lack of reflexes in the distribution of *bətaw ‘sister (male
speaking)’ and *ñaRa ‘brother (female speaking)’ in most western Indonesian
languages, or in any known \tsc{Celebic} or Philippine language raises questions
that are discussed in \crf{ch:MorLes}.

\tsc{Taliabu} languages (\crf{ch:Tal}) and \tsc{Sula-Buru} languages
(\crf{ch:SB}) have “retained” the four-term system of basic sibling
terms, as shown in \trf{07-03}.

\begin{table}\tcp{Four-term sibling sets in \tsc{Taliabu} and \tsc{Sula-Buru}}{07-03}
	\begin{threeparttable}\st{6pt}
	\begin{tabular}{lllll}\lsptoprule
local gl.	&	eSib (same sex)	&	ySib (same sex)	&	sister (m.s.)	&	brother (f.s.)\\
PMP	&	*kakak	&	*huaji	&	*bətaw	&	*ñaRa\\
%\midrule
%\mc{5}{c}{\tsc{Taliabu}}\\
\midrule
\ili{Taliabu}	&		&		&	{\hr}foto	&	{\hr}naha\\
\ili{Soboyo}	&	{\hr}kaka	&	{\hr}uli	&	{\hr}foto	&	{\hr}naha\\
\ili{Kadai}	&	{\hr}kaka	&	{\hr}uli	&	{\hr}foto	&	{\hr}naha\\
%\midrule
%\mc{5}{c}{\tsc{Sula-Buru}}\\
\midrule
\ili{Mangole}	&		&		&	{\hr}feta	&	{\hr}nai,\su{†} \cg{(sanasi)}\\
\ili{Sanana}	&	{\hr}kàk	&		&	{\hr}feta	&	{\hr}nai,\su{†} \cg{(sanohi)}\\
\ili{Lisela}	&	{\hr}kaka	&	{\hr}wai	&	{\hr}feta	&	{\hr}naha-t\\
\ili{Buru}	&	{\hr}kai	&	{\hr}wai	&	{\hr}feta	&	{\hr}naha, naha-t\\
\ili{Ambelau}	&	{\hr}ʔai-o	&	{\hr}βai-no	&	{\hr}bera	&	{\hr}naha-ni\\
		\lspbottomrule
	\end{tabular}
		\begin{tablenotes}\tnsize
			\item[†] {
							Final \it{i{\#}} in \ili{Mangole}
							and \ili{Sanana} is unexplained,
							unless it is reduced from \it{naa-i},
							with regular loss of *R in some dialects.
							The alternate words for ‘brother’ in parentheses are assumed
							to be loans, or based on other roots
							(see discussion in \citealt[503]{Bloyd2020}).}
		\end{tablenotes}
	\end{threeparttable}
\end{table}

In contrast to the 4-term set in \tsc{Sula-Buru}
and \tsc{Taliabu} languages, \tsc{\ili{Ambon}-Seram} languages (\crf{ch:AS})
have innovated a 3-term set of sibling terms.
In \tsc{Alune}, \tsc{Atamanu},
and \tsc{Patakai-Manusela}, reflexes of *bətaw (reconstructed
gloss ‘sister (male speaking)’),
have a broader local meaning of `opposite sex sibling' regardless
of who is speaking.
\ili{Alune} has further reduced to a 2-term set
with `opposite sex sibling' vs. `same sex sibling'.
Languages of the \tsc{\ili{Piru} Bay} branch have innovated a new word
**leku “opposite sex sibling”.
These are illustrated in \trf{07-04}.

\begin{table}[p]\centering\tcp{Three-term sibling sets in \tsc{\ili{Ambon}-Seram} languages}{07-04}
	\begin{threeparttable}\st{5pt}
	\begin{tabular}{lllll}\lsptoprule
local gl.	&	eSib (same sex)	&	ySib (same sex)	&	\mc{2}{l}{opposite sex sibling}\\
PMP	&	*kakak	&	*huaji	&	*bətaw & **leku\\ \midrule
\ili{Manipa}	&	{\hr}haha	&		&	\\
\ili{Kelang}	&	{\hr}haha	&		&	\\
\ili{Luhu}	&	{\hr}kaka (< Mly.)\su{†}	&	{\hr}wali	&&	{\hr\hr}leu {}\\
\ili{Piru}	&		&	{\hr}wali	&&	{\hr\hr}leu {}\\
\ili{Larike}	&	{\hr}aʔa	&	{\hr}wale	&&	{\hr\hr}leko {}\\
\ili{Asilulu}	&	{\hr}ʔaʔa-ni	&	{\hr}wali-mu	&&	{\hr\hr}leku {}\\
\ili{Paulohi}	&	{\hr}waʔa	&	{\hr}waʀi\su{Δ}	&&	{\hr\hr}leʔu {}\\
\ili{Amahai}	&	‹wa›	&	\hp{*w}ari	&&	{\hr\hr}reu {}\\
\ili{Kaibobo}	&	{\hr}waʔa {\tl} aʔa	&	\hp{*w}ari	&&	{\hr}‹leu›\su{‖} {}\\
Waisamu	&	{\hr}waʔa	&	\hp{*w}ari	&&	{\hr}‹leu›\su{‖} {}\\
\ili{Kamarian}	&	{\hr}waʔa	&	{\hr}wari	&&	{\hr}‹reu›\su{‖} {}\\
\ili{Tihulale}	&	{\hr}waʔa	&	{\hr}wari	&&	{\hr}‹leu›\su{‖} {}\\
\ili{Rumakai}	&	{\hr}waʔa	&	{\hr}wari	&&	{\hr}‹reu›\su{‖} {}\\
\ili{Haruku} (\ili{Karu})	&	{\hr}waʔa	&	{\hr}wari	&&	{\hr}‹reu›\su{‖} {}\\
\ili{Haruku} (\ili{Pelauw})	&	\cg{(laha-nu)}	&	{\hr}wari	&&	{\hr}‹leu›\su{‖} {}\\
\ili{Haruku}\su{‡}	&	{\hr}waʔa	&	{\hr}wari	&&	{\hr\hr}reʔu {}\\
\ili{Saparua} (\ili{Hatawano})	&	{\hr}waʔa	&	{\hr}wari	&&	 \\
\ili{Saparua}	&	\cg{(u-wau)}	&	{\hr}wari-u	&&	{\hr\hr}u-leu {}\\
\ili{Nusalaut} (Nalahia)	&	{\hr}waʔa	&		&&	 \\
\ili{Nusalaut}	&	\cg{(waow)}	&	{\hr}wari	&&	{\hr}‹leu›\su{‖} {}\\
N. \ili{Alune}	&		&	{\hr}kwali	&	{\hr}beta-i&\\
W. \ili{Alune}	&		&		&	{\hr}weta, &{\hr}‹deu›\su{‖} {}\\
\ili{Yalahatan}	&	{\hr}ʔaʔa-i	&	{\hr}wali-n	&	{\hr}heta-i&\\
\ili{Atamanu}	&	{\hr}aʔa	&	{\hr}wali	&	{\hr}feta&\\
S. \ili{Nuaulu}	&	{\hr}kaka-i	&	{\hr}wani-ni	&	{\hr}hota-i&\\
N. \ili{Nuaulu}	&	{\hr}kaka-i	&	{\hr}wani-n	&	{\hr}fota-i&\\
\ili{Manusela}	&		&		&	{\hr}hota-mu\\
		\lspbottomrule
	\end{tabular}
		\begin{tablenotes}\tnsize
			\item[†]	{The expected \ili{Luhu} reflex with glottal is found in the compound
								\it{wali-ʔaʔa} ‘same generation relatives’ \citep[145, note 112]{Collins1983}.}
			\item[‡]	{\ili{Haruku} alternate data from \citet[137--151, {[}list 23{]}]{Stokhof1982a}.}
			\item[Δ]	{\ili{Paulohi} ‹r› is explicitly described as uvular
								(\citealp[4]{Stresemann1918}, \citealp[24]{Stresemann1927}).}
			\item[‖] {Words from \citet[334]{vanEkris1865} possibly have a medial
								glottal stop which is not represented.}
		\end{tablenotes}
	\end{threeparttable}
\end{table}

Note that \ili{Alune} has reduced to a 2-term set: \it{kwali} ‘same
sex sibling’ and \it{beta{\nbhyp}i} ‘opposite sex sibling’.
\ili{As} with many languages in central Maluku,
these terms are modified by spatial terms ‘in front (= older)’
and ‘in back, behind (= younger)’ (the latter usually from reflexes
of PMP *ma-udəhi ‘behind, last’).
So \ili{Alune} \it{kwali mena} ‘older (in front) same sex sibling’
and \it{kwali muli} ‘younger (behind) same sex sibling’;
\it{beta mena} ‘older opposite sex sibling’
and \it{beta muli} ‘younger opposite sex sibling’.
We revisit the issue of sibling terms in \srf{15-04-02-04},
where we look at the distribution of opposite sex sibling terms
across Wallacea and beyond.
